\section{Introduction}
\label{sec:intro}
Many people who pay for a gym membership use it far less than they had planned \citep{dellavigna2006}. Models of present bias explain this pattern: a person who gives extra weight to immediate costs keeps postponing an activity whose costs come now and whose benefits come later, and she may not foresee that she will keep postponing it \citep{laibson1997,odonoghue1999}. Paying people to exercise raises attendance while the payments last, but the effects often fade once the payments stop \citep{gneezy2012}.

Two cheap tools address the problem from different sides. A commitment contract raises the cost of skipping the gym for a person who expects to be tempted later \citep{bryan2010}. A reminder makes a planned activity salient at the moment of choice, which helps people who are forgetful rather than impatient \citep{karlan2016}. Evidence on each tool is growing, but the two are rarely tested together, even though the case for combining them is intuitive: a contract gives people a reason to go, and a reminder tells them when.

This paper reports a field experiment with 2,412 employees of a large public university who held a membership at the campus recreation center. We randomized their departments to four arms in a two-by-two design: a control group, text-message reminders, the offer of a self-funded commitment contract, and the offer combined with reminders. A contract required at least 12 visits in eight weeks. Participants chose their own stake, which was refunded if they met the target and donated to a food bank if they did not. Attendance comes from the recreation center's entry records for the 12 weeks before and the 20 weeks after the start of the intervention.

Offering a commitment contract caused a large increase in attendance while the contracts were in force. About a third of the employees offered a contract signed one, and the offer raised visits by 0.31 per week. Reminders did little on their own but made the offer more effective. After the contracts expired, the effect of the offer on visits fell to about a third of its size.

The paper makes three contributions. First, it tests whether reminders and commitment contracts are complements, in a population that has already taken the first step of buying a membership. Second, it randomizes departments rather than individuals, which limits spillovers among colleagues who exercise together. Third, it relates the take-up of contracts to stated self-control problems and compares the cost of the contracts with their effects.

Section~\ref{sec:literature} reviews related work. Section~\ref{sec:design} describes the design and Section~\ref{sec:data} the data. Section~\ref{sec:results} reports the results and Section~\ref{sec:robustness} robustness checks. Section~\ref{sec:discussion} discusses mechanisms and costs, and Section~\ref{sec:conclusion} concludes.
